What is the one-dimensional distribution of the stock price in a jump-diffusion
model with constant drift and volatility, a Poisson counting process and i.i.d. 
log-normal multiplicative jump sizes? Explain.

These models include occasional jumps besides the Brownian diffusion term.
The SDE generalizes Geometric Brownian motion:
\( \frac{dS(t)}{S(t-)} = \mu dt + \sigma dW(t) + dJ(t),\ W \text{ Q-BM},\ J \perp W \)
where the summands are the drift, diffusion and jump part respectively. 

Jump process \( J \)
Random arrival times  \( 0 < \tau_1 < \tau_2 < ... \)
Counting process \( N(t) = \sup \{n : \tau_n \leq t\} \)
Jump sizes \( Y_1 - 1, Y_2 - 1, \dots \) for random variables \( Y_1, Y_2, \dots \)
Then:

\( J(t) = \sum_{j=1}^{N(t)} (Y_j - 1) \)

The symbol \( dJ(t) \) denotes the jump at time t, i.e.,

\( dJ(t) = \begin{cases}
Y_j - 1, & \text{if } t = \tau_j \\
0, & \text{if } t \notin \{ \tau_1, \tau_2, ... \}
\end{cases} \)

which in turn implies \( S(\tau_i) = S(\tau_i -) + S(\tau_i -)(Y_i - 1) = S(\tau_i-)Y_i \)

In the case of multiplicative jumps one computes
\( S(t) = S(0) \cdot \exp \left\{ (\mu - \frac{1}{2} \sigma^2) t + \sigma W(t) \right\} \cdot \prod_{j=1}^{N(t)} Y_j \)

Some special cases allow for some simplification: Poisson and lognormals
Counting process \(N\) is a Poisson process with rate \(\lambda\), i.e., interarrival times \(T_{j+1} - T_j\) independent with
\( \mathbb{P}(T_{j+1} - T_j \le t) = 1 - e^{-\lambda t}, \quad t > 0. \)

Jump sizes are governed by random variables \(Y_1, Y_2, ...\) iid \(\stackrel{\text{iid}}{\sim} LN(a, b^2)\), i.e., \(\ln Y_j \sim N(a, b^2)\).
This implies that for any \(n\) we have \(\prod_{j=1}^n Y_j \sim LN(an, b^2n)\)

Distribution of \(S\) Conditional on \(N(t) = n\):
\( S(t) = S(0) \cdot \exp \{(\mu - \frac{\sigma^2}{2}) t + \sigma W(t)\} \cdot \prod_{j=1}^n Y_j \sim S(0) \cdot LN((\mu - \frac{\sigma^2}{2})t, \sigma^2 t) \cdot LN(an, b^2n) = \), \( LN(\ln S(0) + (\mu - \frac{\sigma^2}{2})t + an, \sigma^2 t + b^2n)\) with distribution function \(F_{n,t\)} \)

Since \(N\) is Poisson-distributed, we obtain a Poisson mixture of lognormals:
\( \mathbb{P} (S(t) \le x) = \sum_{n=0}^\infty e^{-\lambda t} \frac{(\lambda t)^n}{n!} F_{n,t}(x) \)


Describe a stock price model under a martingale measure with constant drift and
volatility, a Poisson counting process and iid multiplicative jump sizes. 

(cont. question) How does its SDE look like? (But note that jump diffusion models are usually incomplete
such that infinitely other martingale measures exist. Here, we just fix one!)

We try to identify an equivalent measure \( Q \) such that \((S(t)e^{−rt})_t \) is a \( Q \)-martingale.

\( \frac{dS_t}{S_{t-}} = \mu_1 dt + \sigma dW_t + dJ_t - \mu_2 dt \)

if \( W \) is a \( Q \)-BM then
\( \mu_1 = r \) the risk free interest rate
\( \mu_2 \) needs to compensate for the jumps \( J(t) = \sum_{j=1}^{N(t)} Y_j - 1 \)
and in general it holds that \( \sum_{j=1}^{N(t)} h(Y_j) - \lambda \mathbb{E}[h(Y_1)]t  \) is martingale
so for our case \( h(y) = y - 1 \) we have \( \mu_2 = \lambda \mathbb{E}[h(Y_1)] = \lambda(\mathbb{E}[Y_1] - 1) \)

and we conclude 
\( \frac{dS_t}{S_{t-}} = rdt + \sigma dW_t + dJ_t - (\lambda(\mathbb{E}[Y_1] - 1)dt \)


Specify an exact simulation of jump diffusion model with Poisson counting process
for a fixed time discretization. How can you even include an exact simulation of
jump times and the distribution of the process at these jump times?

The goal is to simulate the process at dates that are fixed a priori \( 0 = t_0 < t_1 < ... < t_n \):
Our assumptions ar e
\( N \) is a Poisson process
\(Y_1, Y_2, ...\) are iid, but we do not require that their distributions are lognormal; instead the common distribution can flexibly be chosen
The processes \( N, W \) and random variables \(\{Y_1, Y_2, ...\}\) are mutually independent.
This leads to the following time discretization of the log process:
\(X(t_{i+1}) = X(t_i) + (\mu - \frac{1}{2} \sigma^2) (t_{i+1} - t_i) + \sigma \cdot [W(t_{i+1}) - W(t_i)] + \sum_{j=N(t_i)+1}^{N(t_{i+1})} \log Y_j\)

The corresponding algorithm is

1. generate \(Z \sim N(0,1)\)
2. generate \(N \sim \text{Poisson}(\lambda(t_{i+1}-t_{i}))\);
   if \(N = 0\), set \(M = 0\) and go to Step 4
3. generate \(\log Y_1, ..., \log Y_N\); set \(M = \log Y_1 + \dots + \log Y_N\)
4. set 
\(X(t_{i+1}) = X(t_{i}) + (\mu - \frac{1}{2} \sigma^2) (t_{i+1} - t_i) + \sigma \sqrt{t_{i+1} - t_i} \cdot Z + M\) 

Remark:
• If \(Y_j \sim LN(a, b^2)\), then
\(\sum_{j=1}^{n} \log Y_{j} \sim an + b \sqrt{n} \mathcal{N}(0,1)\).
This implies that Step 3 can be replaced by Step 3':
3' generate \(Z' \sim N(0, 1)\); set \(M = aN + b \sqrt{N} Z'\)
• A similar strategy can be chosen if \(\ln Y_j\) is gamma distributed.


To also simulate the jump times, observe that
between jumps the process behaves like GBM
and additionally there are jumps at the jump times.

\(S(\tau_{j+1}-) = S(\tau_{j}) \cdot \exp \left\{ (\mu - \frac{1}{2} \sigma^2) (\tau_{j+1} - \tau_j) + \sigma [W(\tau_{j+1}) - W(\tau_j)] \right\}\)

\(S(\tau_{j+1}) = S(\tau_{j+1}-) \cdot Y_{j+1}\)

Setting \(X(t) = \ln S(t)\), we get

\(X(\tau_{j+1}) = X(\tau_{j}) + (\mu - \frac{1}{2} \sigma^2) (\tau_{j+1} - \tau_j) + \sigma [W(\tau_{j+1}) - W(\tau_j)] + \log Y_{j+1}\)

1. Generate \(R_{j+1}\) from an exponential distribution with mean \(\frac{1}{\lambda}\) 
2. Generate \(Z_{j+1} \sim \mathcal{N}(0,1)\)
3. Generate \(\log Y_{j+1}\)
4. Set \(\tau_{j+1} = \tau_{j} + R_{j+1}\) and

\(X(\tau_{j+1}) = X(\tau_{j}) + (\mu - \frac{1}{2} \sigma^2) R_{j+1} + \sigma \sqrt{R_{j+1}} \cdot Z_{j+1} + \log Y_{j+1}\)

To include an intermediate time step \( t \) into an already run simulation one can save computational effort by the observation

Consider the times
\(\tau_j < t < \tau_{j+1}\)
Then we obtain that

\(S(t) = S(\tau_{j}) \cdot \exp \left\{ (\mu - \frac{1}{2} \sigma^2) (t - \tau_j) + \sigma [W(t) - W(\tau_j)] \right\}\)
\(S(\tau_{j+1}) = S(t) \cdot \exp \left\{ (\mu - \frac{1}{2} \sigma^2) (\tau_{j+1} - t) + \sigma [W(\tau_{j+1}) - W(t)] \right\} \cdot Y_{j+1}\)

What is a Levy process? What are infinitesimal distributions? What is a subordinator?

A Levy process is equivalent to a process \( X \) that satisfies any one tfoa:
1. \( X \) has stationary, independent increments and \( X(t) \to X(s) \) in law, when \( t \to s \)
2. \( X \) can be written as \( X = \mu dt + \sigma dW_t + J_t\) where \( W_t, J_t \) are independent
and \( J_t \) is a pure jump Levy process.
3. \( X(1) \sim F \), where \( F \) is an infinitesimal divisible distribution.
and conversely for \( F \) there always exists a Levy process \( Y \) with \( Y(1) \sim F \)

thus BMs, Poisson processes etc. are all special cases of Levy processes.

A random variable Y has an infinitesimal divisible distribution
if and only if for all 
\( n \in \mathbb{N} \) there are \( Y_1 , \dots, Y_n \) i.i.d. such that \( Y = \sum^n_{i=1} Y_i^{(n)}  \).

Pure jump Levy processes are categorized into
Finite activity: There are finitely many jumps on any compact time interval,
which is the same as as compound Poisson process \(  \sum_{i=1}^{N(t)} Y_i\) with \( Y_i \) i.i.d. and \( N ~ \text{Pois}(\lambda) \) 
Infinite activity
There are infinitely many jumps on any compact time interval.
In this case the increments of the process can still be simulated but not all jumps.

Typically, one assume for a Lévy process \( X \) that \( X(0) = 0 \).
A Lévy process \( X \) is called a subordinator, if \( X \) is increasing. This corresponds precisely to positive infinitesimal distributions, i.e.,
\( X(1) \geq 0 \).
An interesting and tractable class of Lévy processes is constructed as time changes of Brownian motion using subordinators:
In this case, one takes a Brownian motion \( W \) and a subordinator \( G \) with \( W \bot G \).
Then \( (W(G(t))_t \geq 0 \) is again a Lévy process.


Define a variance gamma process, and provide a simulation algorithm in the special
case of only three parameters.

The gamma distribution is infinitely divisible and positive.
This implies that that each gamma distribution corresponds to a subordinator, a gamma process.

Simulating a gamma process
• Let \(X\) be a Lévy subordinator with \(X(1) \sim \text{Gamma}(\alpha, \beta)\).
• Consider a time grid \(0 = t_0 < t_1 < ... < t_n\).
• On the time grid, an exact simulation of the process requires simply an exact simulation of the increments
\(X(t_{i+1}) - X(t_i) \sim \text{Gamma}(\alpha \cdot (t_{i+1} - t_i), \beta)\)

A gamma process \(G\) is a subordinator that can be used to define a time change of a Brownian motion with drift \(W\):
\(VG_{t}=W(G(t)).\)
The process \( VG \) is called a variance gamma process.
One can show that the variance gamma process is in fact the difference of two gamma processes.

Simulation of variance gamma process:

We focus on the following special case with only three parameters:
\( W \) has drift \(\mu\) and volatility \(\sigma\)
\(G(1) \sim \text{Gamma}(\frac{1}{\beta}, \beta)\)
We consider again a time grid \(0 = t_0 < t_1 < ... < t_n\)
Algorithm:
1. generate \(Y \sim \text{Gamma}((t_{i+1} - t_i)/\beta, \beta)\)
2. generate \(Z \sim \mathcal{N}(0,1)\)
3. set \(X(t_{i+1}) = X(t_{i}) + \mu Y + \sigma \sqrt{Y} \cdot Z\)

State in the context of one-dimensional random variables the Monte Carlo estima-
tion problem, the crude method and the approach using control variates.

The flash card also answers the questions
Compute the optimal coefficient and optimal variance.
How can the optimal coefficient be estimated using simulated data?
How does this method generalize to multiple controls?


Problem: Estimate \(\mathbb{E}[Y]\) by simulation!
Simulation: Available: samples \((X_1, Y_1), (X_2, Y_2), ...\) iid such that \(Y_i \sim Y, X_i \sim X, \mathbb{E}[X]\) known
a) Naive approach:
\(\hat{Y}_n := \frac{1}{n} \sum_{i=1}^n Y_i \rightarrow \mathbb{E}[Y]\) a.s. and in \(L^1\) by LLN
b) Alternative approach: Define \(Y_i(b) := Y_i - b(X_i - \mathbb{E}[X])\)
\(\Rightarrow \hat{Y}_n(b) := \frac{1}{n} \sum_{i=1}^n Y_i(b) \rightarrow \mathbb{E}[Y]\) a.s. and in \(L^1\) by LLN
Observe: \(\text{Var}(\hat{Y}_n) = \frac{1}{n} \text{Var}(Y_i(b))\), if \(X, Y \in L^2\)
\( \sigma^2(b) := \text{Var}(Y_i(b)) = \sigma_Y^2 - 2b \sigma_X \sigma_Y \rho_{XY} + b^2 \sigma_X^2\) with \(\sigma_X^2 = \text{Var}(X), \sigma_Y^2 = \text{Var}(Y), \rho_{XY} = \frac{\text{Cov}(X, Y)}{\sigma_X \sigma_Y} \)

Since \(\hat{Y}_n(b)\) is unbiased, smallest MSE is the same as smallest variance.
Minimizing quadratic function \(b \mapsto \sigma^2(b)\) gives optimal coefficient and optimal variance
\(b^{*} = \frac{\text{Cov}(X, Y)}{\text{Var}(X)}, \ \sigma^2(b^{*}) = (1 - \rho_{XY}^2) \sigma_Y^2 (0)\)
\(\rightsquigarrow\) Correlation determines efficiency of method.

Estimating \(b^{*}\): Typically, \(\sigma_Y, \rho_{XY}\) unknown. Use regression estimate:
\( \hat{b}_n = \frac{\sum_{i=1}^n (X_i - \bar{X}_n)(Y_i - \bar{Y}_n)}{\sum_{i=1}^n (X_i - \bar{X}_n)^2} \)

Method can be generalized to the multidimensional case
\(Y\) random variable, \(X = (X^{(1)}, \dots, X^{(d)})^T\) with covariance matrix
\(\begin{pmatrix}
\Sigma_X & \Sigma_{XY} \\
\Sigma_{XY}^T & \sigma_Y^2 
\end{pmatrix}\) 
which is in \(\mathbb{R}^{(d+1) \times (d+1)}\)
Apply \(\bar{Y}_n(b) = \bar{Y}_n - b^T (\bar{X}_n - \mathbb{E}[X])\)
for \(b^{*} = \Sigma_X^{-1} \Sigma_{XY} \in \mathbb{R}^d\).

Missing the answer to "Discuss possible applications, i.e., “underlying assets” and “tractable dynamics”."

Explain the key ideas of antithetic variates. What is the approximate effort? Show
that antithetic variates reduce variance.

Goal: Estimate \(\mathbb{E}[Y]\) for \(Y = f(Z_1, \dots, Z_d)\) with \(f: \mathbb{R}^d \rightarrow \mathbb{R}\), \(Z = (Z_1, \dots, Z_d) \sim \mathcal{N}(0, I_d)\)
• Observation:
With \(f_0(z) := \frac{f(z) + f(-z)}{2}\), we get \(\mathbb{E}[f_0(Z)] = \mathbb{E}[Y]\).
Replacing samples of \(Y\) by samples of \(f_0(Z)\) is antithetic sampling.

Variance reduction:
Define \(f_1(z) := \frac{f(z) - f(-z)}{2}\). Then \(f = f_0 + f_1\), where \(f_0\) is the symmetric part and \(f_1\) is the asymmetric part.
Then \(\text{Var}(f_0(Z)) = \text{Var}(Y) - \text{Var}(f_1(Z))\), since
\(\text{Var}(Y) = \text{Var}(f(Z)) = \text{Var}(f_0(Z) + f_1(Z)) = \text{Var}(f_0(Z)) + \text{Var}(f_1(Z)) + 2 \text{Cov}(f_0(Z), f_1(Z))\)
\(\text{Cov}(f_0(Z), f_1(Z)) = \mathbb{E} \left[ \left( \frac{f(Z) + f(-Z)}{2} - \mathbb{E}[Y] \right) \left( \frac{f(Z) - f(-Z)}{2} \right) \right]\)
\(= \frac{1}{4} \mathbb{E} \left[ (f(Z)^2 - f(-Z)^2) \right] - \frac{\mathbb{E}[Y]}{2} [\mathbb{E}[f(Z)] - \mathbb{E}[f(-Z)]] = 0\)

Antithetic sampling removes the variance of the antisymmetric part.
This idea generalizes to other measure-preserving transformations: systematic sampling.

This also answers the questions
In the context of Gaussians instead of uniform simulation factors, define the addi-
tive decomposition of a function into a symmetric and antisymmetric part.
Prove that antithetic sampling removes the variance of the antisymmetric part.

Idea: negative dependence between pairs of replication
If \(U \sim \text{unif}[0, 1]\), then also \(1-U \sim \text{unif}[0, 1]\). If \(F\) is a distribution function, then \(F^{-1}(U) \sim F, F^{-1}(1-U) \sim F\).

Example: \(Z \sim \mathcal{N}(0, 1), U = \Phi(Z) \Rightarrow Z = \Phi^{-1}(U) \sim \Phi^{-1}(1-U) \sim \mathcal{N}(0, 1)\)
Problem: Given \(Y \sim F\). Compute \(\mathbb{E}[Y]\) by simulation!
Method:

Generate \(U_1, U_2, ... \sim \text{unif}[0, 1]\) iid
Construct pairs \((Y_1, \tilde{Y}_1), (Y_2, \tilde{Y}_2), ...\) by setting \(Y_i = F^{-1}(U_i), \tilde{Y}_i = F^{-1}(1-U_i)\), \(i = 1, 2, ...\)
Antithetic variates estimation:

\(\hat{Y}_{\text{AV}, n} = \frac{1}{2n} \left( \sum_{i=1}^n Y_i + \sum_{i=1}^n \tilde{Y}_i \right) = \frac{1}{n} \sum_{i=1}^n \frac{Y_i + \tilde{Y}_i}{2} \rightarrow \mathbb{E}[Y]\) by LLN

(a.s., in \(L^1\), in \(L^2\) if everything in \(L^2\))


Explain the basic idea of stratified sampling. How are general stratification vari-
ables and general proportions integrated?

General stratification variable:
\(X \ \mathbb{R}^d\)-valued random variable, \(A_1, A_2, \dots, A_k\) partition of \(\mathbb{R}^d\), \(p_i = \mathbb{P}[X \in A_i]\)
then \(\mathbb{E}[Y] = \sum_{i=1}^k p_i \mathbb{E}[Y | X \in A_i]\)
For each \(i\), sample from \(\mathcal{L}(Y | X \in A_i)\).

General proportions \(n_1, \dots, n_k\):
Set \(q_i := \frac{n_i}{n}\), then we obtain the estimator
\(\hat{Y} = \frac{1}{n} \sum_{i=1}^k \frac{p_i}{q_i} \sum_{j=1}^{n_i} Y_{ij}\)


In particular for a given cumulative distribution function \(F: \mathbb{R} \rightarrow [0, 1]\) with quantile function \(F^{-1}(u) = \inf \{ x : F(x) \geq u \}\), 
\(u \in [0, 1]\)
Let \(p_1, p_2, ..., p_k > 0\) with \(\sum_{i=1}^{k} p_i = 1\) and set

\(a_1 = F^{-1}(p_1)\)
\(a_2 = F^{-1}(p_1 + p_2), ..., a_k = F^{-1}(1)\)

With \(a_0 = -\infty\) define strata

\(A_1 = (a_0, a_1], A_2 = (a_1, a_2], ..., A_k = \begin{cases}
(a_{k-1}, a_k], & \text{if } a_k < \infty \\
(a_{k-1}, \infty), & \text{otherwise}
\end{cases}\)

Then \(\mathbb{P}[Y \in A_i] = p_i\).
If \(U \sim \text{unif}[0, 1]\), then \(F^{-1}(V) \sim \mathcal{L}(Y | Y \in A_i)\) with \(V = p_1 + \dots + p_{i-1} + U \cdot p_i\).
The empirical distribution of stratified samples provides a better approximation of the original distribution than the empirical distribution of iid-samples of the original distribution.

If we denote by \( V_i \) the event \( Y|X \in A_i \) then the special case of \( F = \text{unif}[0, 1] \) is given by
\( V_i = \frac{i-1}{n} + \frac{U_i}{n}\), where \( U_i \sim \text{unif}[0, 1]\) and sampled once for each \( A_i \).

Estimate: \(\mu := \mathbb{E}[Y]\)
Strata: \(\{X \in A_i\}, i = 1, \dots, k \rightarrow n_i\) samples \(Y_{ij}\) of 
\(\mathcal{L}(Y|X \in A_i), i = 1, \dots, k,\ n = \sum_{i=1}^k n_i,\ q_i = \frac{n_i}{n},\ \mathbb{P}[X \in A_i] = p_i\)

Estimator: \(\hat{Y} = \sum_{i=1}^k p_i \sum_{j=1}^{n_i} Y_{ij} = \frac{1}{n} \sum_{i=1}^k \frac{p_i}{q_i} \sum_{j=1}^{n_i} Y_{ij}\)
\(\mathbb{E}[\hat{Y}] = \mu\)

If \(n\) is increased such that \(q_i \sim\) const., then convergence to \(\mu\).

Variance:
Set \(\mu_i := \mathbb{E}[Y_{ij}] = \mathbb{E}[Y | X \in A_i]\)
Then \(\sigma_i^2 := \text{Var}(Y_{ij}) = \text{Var}(Y | X \in A_i)\)
Then 
\(\text{Var}(\hat{Y}) = \sum_{i=1}^k p_i^2 \text{Var} \left( \frac{1}{n_i} \sum_{j=1}^{n_i} Y_{ij} \right)\) 
\(= \sum_{i=1}^k \frac{p_i^2 \sigma_i^2}{n_i} = \frac{1}{n} \sum_{i=1}^k \frac{p_i^2 \sigma_i^2}{q_i} \)
\(= \frac{\sigma^2(q)}{n}\) with \(\sigma^2(q) = \sum_{i=1}^k \frac{p_i^2}{q_i}\sigma_i^2  \)

Using the example of the uniform distribution on a high-dimensional unit hyper-
cube, explain the curse of dimensionality.

Stratification creates problems in higher dimensions.
Consider a d-dimensional hypercube \([0, 1]^d\) and a uniform distribution on \([0, 1]^d\).
Method:
j-th coordinate \(\rightarrow K_j\) intervals of equal length
Form of each stratum: \(\prod_{i=1}^d \left[ \frac{j_i - 1}{K_i}, \frac{j_i}{K_i} \right) =: A\) with probability \(\frac{1}{K_1 \dotsm K_d}\)
If \(Y \sim \text{unif}[0, 1]^d\), then \(\mathcal{L}(Y | Y \in A) = \text{unif } A\).

Sampling: \(U_1, \dots, U_d \sim \text{unif}[0, 1]\), \(V_j = \frac{j_i - 1 + U_j}{K_j}, j = 1, \dots, d \Rightarrow V \sim \text{unif } A\)
Problem: \(K_1 \dotsm K_d\) strata \(\rightarrow\) too many in high dimensions (\(d > 5\))


Optimize the constant of the variance for stratified r.v. by choosing optimal proportions. 
This allocation depends on probabilities of strata and the variances of the samples on the strata. 
Discuss this.

We know
\( \text{Var}(\hat{Y}) = \frac{\sigma^{2}(q)}{n} \)
and \( \sigma^{2}(q) = \frac{p_i^2}{q_i}\sigma_i^2 \).


Optimization leads to \(q_i^* = \frac{p_i \sigma_i}{\sum_{l=1}^k p_l \sigma_l}, i = 1, ..., k,\) with \(\sigma^2(q^*) = \left( \sum_{l=1}^k p_l \sigma_l \right)^2\) 
\(\sigma^2(p) - \sigma^2(q^*) = \sum_{i=1}^k p_i \sigma_i^2 - \left( \sum_{i=1}^k p_i \sigma_i \right)^2 = \sum_{i=1}^k p_i (\sigma_i - \sum_{j=1}^k p_j \sigma_j)^2\)
Results: When the probability-weight fluctuations of the standard deviations \(\sigma_i\) across the strata become large, the improvement when choosing the optimal 
\(q^*\) is also larger.

Computation cost: possible not uniform across strata \(\Rightarrow\) efficiency requires a more complete analysis 

Explain the stratification of Brownian motion based on its terminal value.

Or more generally this methods is coined terminal stratification:
Motivation: stochastic price process and option that possibly depends on the whole path, 
but terminal price is a good stratification variable

Example: Stratify Brownian Motion based on its terminal value
time points \(0 = t_0 < t_1 < ... < t_m = \) terminal time
sample \(W(t_1), W(t_2), ..., W(t_m)\)
   stratify \(W(t_m)\) using the inverse transform method
   generate intermediate values \(W(t_1), W(t_2), ..., W(t_m)\) by Brownian bridge sampling

Specific example: k equiprobable strata and a proportional allocation of samples
\(U_1, U_2, ..., U_k \sim \text{unif}[0, 1]\) iid, \(V_i = \frac{i-1}{k} + \frac{U_i}{k}\)
\(\Rightarrow \sqrt{t_m} \Phi^{-1}(V_1), ..., \sqrt{t_m} \Phi^{-1}(V_k)\) stratified samples of \(W(t_m)\)
   gaps are filled by Brownian bridge construction:
\(j < k < l: \mathcal{L}(W(t_k) | W(t_j), W(t_l)) \sim\)
\(\mathcal{N} \left( \frac{(t_l - t_k) W(t_j) + (t_k - t_j) W(t_l)}{t_l - t_j},\ \frac{(t_l - t_k)(t_k - t_j)}{t_l - t_j} \right) \)

Aim: Estimate \(\mathbb{E}[Y]\) for some random variable Y
Available: \((X_1, Y_1), (X_2, Y_2), \dots\) iid
Standard estimator: \(\bar{Y}_n = \frac{1}{n} \sum_{i=1}^n Y_i \rightarrow \mathbb{E}[Y]\) for \(n \uparrow \infty\) by LLN

What is poststratification?

Poststratified estimator
Define:
\(N_{i,n} = \sum_{j=1}^n 1_{\{X_j \in A_i\}}, \ S_{i,n} = \sum_{j=1}^n 1_{\{X_j \in A_i\}} Y_j\)
Then: \(\hat{Y} = \sum_{i=1}^k \frac{N_{i,n}}{n} \frac{S_{i,n}}{N_{i,n}}\) with \(\frac{N_{i,n}}{n} \rightarrow p_i = \mathbb{P}[X \in A_i]\) as \(n \uparrow \infty\)
Suppose that \(p_i, i = 1, \dots, k\), are known or can be easily estimated, then the poststratified estimator 
\(\hat{Y} = \sum_{i=1}^k p_i \frac{S_{i,n}}{N_{i,n}}\) 
has a reduced variance.
